\documentclass{resume}

\begin{document}

\cvheader
  {Rodrigo Roque}
  {Engenheiro de Software Sênior · Staff / Especialista}
  {São Paulo, SP · Brasil\cvsep
   \href{mailto:rodrigogoncalvesroque@gmail.com}{rodrigogoncalvesroque@gmail.com}\\
   \href{https://github.com/roquerodrigo}{github.com/roquerodrigo}\cvsep
   \href{https://linkedin.com/in/roque-rodrigo}{linkedin.com/in/roque-rodrigo}}

\cvsection{Resumo}

Engenheiro de software sênior com mais de 10 anos em backend e \textbf{sistemas
distribuídos}, atualmente no \textbf{Itaú Unibanco} (Soluções de Investimentos). Atuo
como referência técnica em sistemas críticos e de alto volume: defino \textbf{arquitetura
orientada a eventos}, elevo a confiabilidade com práticas de \textbf{SRE} e multiplico o
time por meio de mentoria, code review e disseminação de padrões na comunidade de
engenharia. Minha trajetória passa por fintech de pagamentos, SaaS e o maior banco
privado da América Latina, com \textbf{Java}, \textbf{Python} e \textbf{AWS}. Montei um
time de engenharia do zero, liderei squads e reduzi 35\% de custo de cloud. Busco atuar
como \textbf{Staff Engineer / Especialista}, direcionando decisões técnicas de amplo impacto.

\cvsection{Experiências Profissionais}

\cvevent{Engenheiro de Software Sênior}{Abr 2025 – Presente}
        {Itaú Unibanco}{São Paulo, SP · Híbrido}
\cvsummary{Desenvolvo soluções para o domínio de Perfil do Investidor / Suitability
na comunidade de Soluções de Investimentos, com Java e Python em arquitetura
orientada a eventos na AWS.}
\begin{cvitems}
  \item \textbf{Modernização:} migração de sistemas legados para a AWS, evoluindo a
        arquitetura dos serviços com IaC (Terraform), containerização (Docker/Kubernetes)
        e pipelines CI/CD (GitHub Actions).
  \item \textbf{Confiabilidade:} práticas de SRE com observabilidade de ponta a ponta
        via Datadog (logs, métricas, traces), SLOs, resiliência e resposta a incidentes.
  \item \textbf{Integrações:} desenho e implementação de integrações orientadas a
        eventos com Kafka e APIs internas.
  \item \textbf{IA:} estudos e PoCs de IA aplicada ao negócio e ao ciclo de
        desenvolvimento, além de fomento a ferramentas de IA na produtividade do time.
  \item \textbf{Stakeholders:} tradução de requisitos de negócio em soluções técnicas
        viáveis e disseminação de conhecimento na comunidade de engenharia.
\end{cvitems}

\cvevent{Tech Lead}{Set 2024 – Abr 2025}
        {Itaú Unibanco}{São Paulo, SP · Híbrido}
\cvsummary{Liderei tecnicamente a Squad Revisão de Carteira na comunidade de Soluções
de Investimentos, desenvolvendo com Java e Python em infraestrutura AWS.}
\begin{cvitems}
  \item \textbf{Liderança técnica:} decisões de arquitetura, code review e mentoria do
        time, garantindo qualidade e padronização das entregas.
  \item \textbf{Negócio e engenharia:} atuação como ponte no refinamento e priorização
        técnica das demandas da squad.
  \item \textbf{Observabilidade:} implementação de práticas de monitoramento com Datadog
        para os serviços em produção.
\end{cvitems}

\cvevent{Tech Lead}{Mar 2024 – Set 2024}
        {EntrePay}{São Paulo, SP · Híbrido}
\cvsummary{Construí o time de engenharia do zero e estabeleci a cultura técnica da
empresa, definindo processos de code review, CI/CD, testes automatizados e observabilidade.}
\begin{cvitems}
  \item \textbf{Redução de custos:} diminuição de 35\% no custo de cloud ao otimizar a
        alocação de recursos de infraestrutura.
  \item \textbf{Cultura open source:} implantação de stack de observabilidade e DevOps
        (GitLab, Grafana, Prometheus, OpenTelemetry, Sentry), eliminando custos de licenciamento.
  \item \textbf{Documentação técnica:} organização dos projetos com diagramas e ADRs,
        reduzindo o tempo de onboarding de novos desenvolvedores.
  \item \textbf{Pessoas:} ciclos de 1:1, retrospectivas, plano de carreira técnica e
        workshops de disseminação de conhecimento.
\end{cvitems}

\cvevent{Engenheiro de Software Sênior}{Jul 2021 – Dez 2023}
        {Arquivei}{São Carlos, SP · Remoto}
\cvsummary{Atuei na entrega ágil de MVPs para validação de hipóteses de negócio, além
de exercer liderança técnica como referência do time de desenvolvimento.}
\begin{cvitems}
  \item \textbf{Pagamentos PIX:} desenvolvimento de API Node.js/TypeScript de pagamentos
        via PIX com Open Finance, processando uma base de \textasciitilde2 bilhões de
        documentos fiscais para gerar insights de fluxo de caixa.
  \item \textbf{Antecipação de Recebíveis:} da arquitetura à validação — API REST em PHP
        (Laravel) com OpenAPI, BFF GraphQL (Apollo), frontend React e infraestrutura em
        Kubernetes/GCP.
  \item \textbf{Observabilidade:} monitoramento com Prometheus e Grafana, aumentando a
        visibilidade sobre a saúde dos serviços em produção.
  \item \textbf{Compartilhamento de conhecimento:} apresentação de novas tecnologias,
        padrões de projeto e arquitetura (DDD, arquitetura hexagonal) ao time de engenharia.
\end{cvitems}

\cvevent{Desenvolvedor Backend Sênior}{Abr 2020 – Jun 2021}
        {Talentify}{Windermere, FL / Barueri, SP · Remoto}
\cvsummary{Desenvolvi integrações da engine de captura de vagas — principal fonte de
receita da empresa — em ambiente 100\% remoto e internacional.}
\begin{cvitems}
  \item \textbf{Agregador de vagas:} integrações via APIs e web crawlers que mantinham
        \textasciitilde1 milhão de vagas de emprego ativas simultaneamente na base de dados.
  \item \textbf{Engenharia:} microsserviços PHP com arquitetura hexagonal, TDD e
        princípios SOLID, garantindo confiabilidade do fluxo crítico de receita.
  \item \textbf{Práticas:} arquitetura orientada a eventos, code review, CI/CD e Docker
        em pipeline de entrega contínua.
\end{cvitems}

\cvevent{Desenvolvedor Backend Pleno}{Ago 2019 – Abr 2020}
        {Lumanai}{São Paulo, SP · Híbrido}
\cvsummary{Desenvolvi produtos em equipe ágil (Scrum), com refinamento de histórias e
colaboração em times maiores.}
\begin{cvitems}
  \item \textbf{Plataforma de intercâmbio:} migração de aplicação legada de CRM para nova
        arquitetura (API PHP/Laravel + frontend Vue.js), sem interrupção para os clientes.
  \item \textbf{Sistema financeiro:} manutenção e evolução do sistema de gestão financeira
        da empresa em Python/Django.
\end{cvitems}

\cvevent{Desenvolvedor Fullstack (Estágio)}{Jul 2017 – Mar 2019}
        {d'hire — Executive Search Marketplace}{Rio de Janeiro, RJ · Remoto}
\begin{cvitems}
  \item Desenvolvimento de APIs em PHP/Laravel e interfaces web em Vue.js para marketplace
        de executive search.
  \item Sustentação de aplicações Java (Spring MVC) e Angular.js em produção.
  \item Deploys na nuvem utilizando os serviços da AWS.
\end{cvitems}

\cvevent{Estagiário de Suporte}{Mai 2016 – Mai 2017}
        {UNIFEI — Universidade Federal de Itajubá}{Itajubá, MG · Presencial}
\begin{cvitems}
  \item Implementação e customização de funcionalidades da plataforma Moodle e do site
        institucional (Joomla) em PHP.
  \item Suporte aos funcionários do núcleo e aos alunos usuários da plataforma.
\end{cvitems}

\cvevent{Estagiário de Desenvolvimento}{Abr 2015 – Abr 2016}
        {CVS Sistemas}{Itajubá, MG · Presencial}
\begin{cvitems}
  \item Desenvolvimento e manutenção de aplicações web com PHP, HTML, CSS e JavaScript (jQuery).
  \item Manutenção de redes de computadores e suporte aos usuários.
\end{cvitems}

\cvsection{Habilidades}

\textbf{Arquitetura de Software} (DDD, hexagonal, microsserviços, orientada a eventos) ·
\textbf{Sistemas Distribuídos} · \textbf{Java} · \textbf{Python} · \textbf{AWS} ·
\textbf{SRE} · Observabilidade (Datadog, Grafana, Prometheus, OpenTelemetry) · \textbf{Kafka} ·
Terraform · Docker · Kubernetes · GitHub Actions · CI/CD · Spring · APIs REST · GraphQL ·
SQL · Git · Linux · IA aplicada · Liderança Técnica · Mentoria · Metodologias Ágeis (Scrum) ·
Node.js · TypeScript · PHP

\cvsection{Formação Acadêmica}

\cvline{Bacharelado em Ciência da Computação}{2015 – 2019}
\noindent{\small\color{muted}UNIFEI — Universidade Federal de Itajubá}\par
\cvline{Técnico em Informática}{2013 – 2015}
\noindent{\small\color{muted}Senac Minas}\par
\cvline{Técnico em Redes de Computadores}{2013 – 2014}
\noindent{\small\color{muted}Senac Minas}\par

\cvsection{Idiomas}

\cvline{Português}{Nativo}
\cvline{Inglês}{Intermediário — boa leitura e escrita técnica}

\end{document}
